% Part of the TAPL project, under the Apache License v2.0 with LLVM
% Exceptions. See /LICENSE for license information.
% SPDX-License-Identifier: Apache-2.0 WITH LLVM-exception

\documentclass[11pt,letterpaper]{article}
\usepackage[margin=1in]{geometry}
\usepackage{amsmath,amssymb,mathtools}
\usepackage{mathpartir}
\usepackage{xcolor}
\usepackage[normalem]{ulem}
\usepackage{enumitem}
\usepackage[hidelinks]{hyperref}

%% ── Macros (each notation is introduced in the text where it is first used) ──
\newcommand{\step}{\longrightarrow}             % one step
\newcommand{\wstep}{\longrightarrow_{\mathrm{w}}} % one step of the evaluation strategy
\newcommand{\lam}[2]{\lambda #1.\,#2}           % λx. t
\newcommand{\cl}[2]{#1{\sim}#2}                 % t~F, a clad term
\newcommand{\subst}[3]{[#1 \mapsto #2]\,#3}     % [x ↦ s] t
\newcommand{\fn}[1]{\mathit{#1}}                % meta function
\newcommand{\rn}[1]{\textsc{#1}}                % rule name in text
\newcommand{\kw}[1]{\mathtt{#1}}                % keyword in math
\newcommand{\dup}{\kw{dup}}
\newcommand{\drop}{\kw{drop}}
\newcommand{\tm}[1]{\fn{term}(#1)}              % term half
\newcommand{\fh}[1]{\fn{form}(#1)}              % form half

\mprset{sep=1.5em}

%% Numbered points: rules, formula rows and bullets share one counter per section,
%% shown as 1), 2), ... (the section number is not repeated) so each can be referred to.
\newcounter{pt}[section]
\renewcommand{\thesubsection}{\thesection\alph{subsection}}  % 4a, 4b, ...
\renewcommand{\thept}{\arabic{pt})}
\newcommand{\pt}{\refstepcounter{pt}\textup{\thept}}  % number a formula row
\newlist{points}{enumerate}{1}
\setlist[points]{label=\arabic*), leftmargin=*,
  start=\numexpr\value{pt}+1\relax, after={\setcounter{pt}{\value{pointsi}}}}

%% Review marks; removed once the user accepts a change.
\newcommand{\new}[1]{{\color{blue}#1}}          % new or changed text
\newcommand{\old}[1]{{\color{red}\sout{#1}}}    % removed text

%% One rule per line, centered on the page, its name right-aligned in a column.
%% The left column (the rule's number) is as wide as the name column, so the rule is
%% truly centered.
%% Each rule is its own row, so a long list can break across pages.
%% \axiomrow{Name}{conclusion}            a rule with no premises, no bar
%% \rulerow{Name}{premises}{conclusion}   premises over a bar
\newlength{\rulenamewidth}
\setlength{\rulenamewidth}{6.5em}
\newenvironment{rules}{\par\smallskip}{\par}
\newcommand{\rulename}[1]{\makebox[\rulenamewidth][r]{\rn{#1}}}
\newcommand{\ruleline}[2]
  {\par\noindent
   \begin{tabular*}{\linewidth}{@{}l@{\extracolsep{\fill}}c@{\extracolsep{\fill}}r@{}}
   \makebox[\rulenamewidth][l]{\pt} & #2 & \rulename{#1}
   \end{tabular*}\par}
\newcommand{\axiomrow}[2]{\ruleline{#1}{$#2$}\vspace{8pt}}
\newcommand{\rulerow}[3]{\ruleline{#1}{$\inferrule{#2}{#3}$}\vspace{12pt}}

\title{Oymo Semantics\\[8pt]\normalsize\textsc{Draft --- under review}}
\author{}
\date{}

\begin{document}

\maketitle

%% ── Syntax ─────────────────────────────────────────────────────────────

\section{Syntax}\label{sec:syntax}

\[
\begin{array}{l@{\quad}rcll}
 & t, s, F, G & ::= & & \textit{terms} \\
\pt &   & \mid & x & \textit{variable} \\
\pt &   & \mid & \lam{x}{t} & \textit{abstraction} \\
\pt &   & \mid & t\ s & \textit{application} \\
\pt &   & \mid & s \Rightarrow t & \textit{signature} \\
\pt &   & \mid & \dup\ s\ t & \textit{copy} \\
\pt &   & \mid & \drop\ s\ t & \textit{discard} \\
\pt &   & \mid & \cl{t}{F} & \textit{clad term}
\end{array}
\]
\begin{points}
\item $x$ ranges over variables. $t$, $s$, $F$ and $G$ range over the same terms;
  $F$ and $G$ are used where a term serves as a form, and $s$ is a companion of
  $t$ where a rule has two terms.
\item A term is \emph{naked} when its root is not ${\sim}$, and \emph{clad} when
  it is $\cl{t}{F}$. In $\cl{t}{F}$, $t$ is the term side and $F$ the form side.
\item A form is similar to a type, but with no typing rules: nothing checks it.
\item A naked term with no ${\sim}$ inside is read as $\cl{t}{t}$, clad by
  itself. The \emph{halves} of a term go through every construct and stop at
  each ${\sim}$:
  \[
  \tm{\cl{t}{F}} = \tm{t},\qquad \fh{\cl{t}{F}} = \tm{F},
  \]
  and for every other construct both halves are taken part by part, e.g.
  $\tm{t\ s} = \tm{t}\ \tm{s}$, $\fh{\lam{x}{t}} = \lam{x}{\fh{t}}$,
  $\tm{x} = \fh{x} = x$. Both halves are naked terms with no ${\sim}$ inside.
\item Terms are equal up to renaming bound variables; substitution
  $\subst{x}{s}{t}$ never captures a variable.
\end{points}

%% ── Computation ────────────────────────────────────────────────────────

\section{Computation (at the root)}\label{sec:computation}

\begin{rules}
\axiomrow{E-Beta}
  {(\lam{x}{t})\ s \;\step\; \subst{x}{s}{t}}
\axiomrow{E-SigApp}
  {(F \Rightarrow G)\ t \;\step\; \drop\ F\ (\drop\ t\ G)}
\rulerow{E-Dup}
  {s \text{ is a closed value}}
  {\dup\ s\ t \;\step\; t\ s\ s}
\rulerow{E-Drop}
  {s \text{ is a closed value}}
  {\drop\ s\ t \;\step\; t}
\end{rules}
\begin{points}
\item A value is $\lam{x}{t}$ or $F \Rightarrow G$ (Section~\ref{sec:classes});
  it is closed when it has no free variable.
\item \rn{E-SigApp}: applying a signature gives its result; the signature's
  parameter and the argument are discarded.
\item \rn{E-Dup} and \rn{E-Drop} only copy or discard a closed value: work that
  is not done yet is never copied or lost, and no variable is copied or lost
  (Section~\ref{sec:linearity}). Otherwise they wait for run time.
\end{points}

\subsection{Clad terms}

\begin{rules}
\rulerow{E-SplitApp}
  {t \text{ or } s \text{ is clad}}
  {t\ s \;\step\; \cl{(\tm{t}\ \tm{s})}{(\fh{t}\ \fh{s})}}
\rulerow{E-SplitAbs}
  {t \text{ is clad}}
  {\lam{x}{t} \;\step\; \cl{(\lam{x}{\tm{t}})}{(\lam{x}{\fh{t}})}}
\rulerow{E-SplitSig}
  {s \text{ or } t \text{ is clad}}
  {s \Rightarrow t \;\step\; \cl{(\tm{s} \Rightarrow \tm{t})}{(\fh{s} \Rightarrow \fh{t})}}
\rulerow{E-SplitDup}
  {s \text{ or } t \text{ is clad}}
  {\dup\ s\ t \;\step\; \cl{(\dup\ \tm{s}\ \tm{t})}{(\dup\ \fh{s}\ \fh{t})}}
\rulerow{E-SplitDrop}
  {s \text{ or } t \text{ is clad}}
  {\drop\ s\ t \;\step\; \cl{(\drop\ \tm{s}\ \tm{t})}{(\drop\ \fh{s}\ \fh{t})}}
\rulerow{E-Reclad}
  {t \text{ or } F \text{ is clad}}
  {\cl{t}{F} \;\step\; \cl{\tm{t}}{\tm{F}}}
\rulerow{E-Collapse}
  {t \text{ has no } {\sim} \text{ inside}}
  {\cl{t}{t} \;\step\; t}
\end{rules}
\begin{points}
\item The split rules move a construct with a clad part into both sides: the
  term side computes with the term halves, the form side with the form halves.
  For example $(\cl{t}{F})\ s \step \cl{(t\ s)}{(F\ s)}$: the result's form is
  the form applied to the argument.
\item The halves go all the way down, so a naked part with a ${\sim}$ inside is
  split correctly too: $(\lam{z}{\cl{z}{F}})\ (\cl{s}{G})$ gives
  $\cl{s}{F}$ both by \rn{E-Beta} and by \rn{E-SplitApp}.
\item \rn{E-Reclad}: a written form replaces the old one, and a form's own form
  is discarded: $\cl{(\cl{t}{F})}{(\cl{G}{H})} \step \cl{t}{G}$. ${\sim}$ does
  not split like the other constructs, which would give a clad term again.
\item \rn{E-Collapse} follows from reading a naked $t$ as $\cl{t}{t}$.
\item None of these rules changes the halves of a term. After them, both sides
  of every clad term are naked.
\end{points}

%% ── Congruence ─────────────────────────────────────────────────────────

\section{Full \texorpdfstring{$\beta$}{beta}}\label{sec:congruence}

\begin{rules}
\rulerow{E-Abs}{t \step t'}{\lam{x}{t} \;\step\; \lam{x}{t'}}
\rulerow{E-App1}{t \step t'}{t\ s \;\step\; t'\ s}
\rulerow{E-App2}{s \step s'}{t\ s \;\step\; t\ s'}
\rulerow{E-Sig1}{s \step s'}{s \Rightarrow t \;\step\; s' \Rightarrow t}
\rulerow{E-Sig2}{t \step t'}{s \Rightarrow t \;\step\; s \Rightarrow t'}
\rulerow{E-Dup1}{s \step s'}{\dup\ s\ t \;\step\; \dup\ s'\ t}
\rulerow{E-Dup2}{t \step t'}{\dup\ s\ t \;\step\; \dup\ s\ t'}
\rulerow{E-Drop1}{s \step s'}{\drop\ s\ t \;\step\; \drop\ s'\ t}
\rulerow{E-Drop2}{t \step t'}{\drop\ s\ t \;\step\; \drop\ s\ t'}
\rulerow{E-Clad1}{t \step t'}{\cl{t}{F} \;\step\; \cl{t'}{F}}
\rulerow{E-Clad2}{F \step F'}{\cl{t}{F} \;\step\; \cl{t}{F'}}
\end{rules}
\begin{points}
\item $\step$ is full $\beta$: a rule of Section~\ref{sec:computation} may fire
  anywhere, under a $\lambda$ and in either side of ${\sim}$ too.
\item $\step$ is expected to be confluent. Proof idea: each half maps a term to a
  naked term, and a step on a term is zero or more steps on each half; the rules
  for clad terms leave both halves unchanged. Confluence then comes from the
  naked calculus (Church--Rosser).
\end{points}

%% ── Linearity ──────────────────────────────────────────────────────────

\section{Linearity}\label{sec:linearity}

\begin{points}
\item A naked term with no ${\sim}$ inside is \emph{linear} when, for every
  $\lam{x}{u}$ in it, $x$ occurs exactly once in $u$, counting occurrences inside
  inner $\lambda$s.
\item Any term is linear when both of its halves are linear: the term side and
  the form side are each checked, as the naked terms they become after
  splitting. For example $\lam{x}{\cl{x}{F}}$ needs $x$ once in $F$ too, since
  its form half is $\lam{x}{F}$.
\item Free variables are not checked again inside a body: each is checked at its
  own binder, or comes from the environment, which is trusted.
\item $\dup$ and $\drop$ are the only ways to copy or discard; they bind nothing.
  For example $\dup\ x\ (\lam{a}{\lam{b}{a\ b}})$ uses $x$ twice.
\item Reduction keeps a term linear: \rn{E-Beta} puts $s$ where the single $x$
  was, \rn{E-Dup} and \rn{E-Drop} only copy or discard closed terms, and the
  rules for clad terms leave both halves unchanged.
\end{points}

%% ── Strategy ───────────────────────────────────────────────────────────

\section{Evaluation Strategy}\label{sec:strategy}

\begin{rules}
\rulerow{W-Root}{t \step t' \text{ by a rule of Section~\ref{sec:computation}}}
  {t \;\wstep\; t'}
\rulerow{W-App1}{t \wstep t'}{t\ s \;\wstep\; t'\ s}
\rulerow{W-Dup1}{s \wstep s'}{\dup\ s\ t \;\wstep\; \dup\ s'\ t}
\rulerow{W-Drop1}{s \wstep s'}{\drop\ s\ t \;\wstep\; \drop\ s'\ t}
\end{rules}
\begin{points}
\item $\wstep$ is weak and in normal order: it never reduces under a $\lambda$,
  inside a signature, or in an argument.
\item $\wstep$ for clad terms is not written yet.
\end{points}

%% ── Term classes ───────────────────────────────────────────────────────

\section{Term Classes}\label{sec:classes}

\[
\begin{array}{l@{\quad}rcll}
\pt & r & ::= & (\lam{x}{t})\ s \mid (F \Rightarrow G)\ t \mid \dup\ v\ t \mid
  \drop\ v\ t & \textit{reducible} \\
 & & \mid & r\ s \mid \dup\ r\ t \mid \drop\ r\ t & \\
\pt & v & ::= & \lam{x}{t} \mid F \Rightarrow G & \textit{value} \\
\pt & \mathit{ne} & ::= & x \mid \mathit{ne}\ t \mid \dup\ \mathit{ne}\ t \mid
  \drop\ \mathit{ne}\ t \mid \dup\ v\ t \mid \drop\ v\ t & \textit{neutral}
\end{array}
\]
\begin{points}
\item The classes are for naked terms without ${\sim}$; the classes for clad
  terms are not written yet.
\item Every such term is in exactly one class: reducible (it takes a $\wstep$
  step), a value, or neutral.
\item In $r$, $\dup\ v\ t$ and $\drop\ v\ t$ need $v$ closed; in $\mathit{ne}$,
  they need $v$ to have a free variable. Being closed is a side condition, not
  grammar.
\item A neutral term waits for a free variable. Once its binder is applied it may
  become reducible.
\item No term is wrong yet, so there is no stuck class.
\end{points}


\end{document}
